\section{Zadanie 2}
\autor{Aleksander Wieczorek}
Policz funkcje tworzące momenty dla rozkładu jednostajnego (na [a, b]), wykładniczego (z parametrem \(\lambda\)) i Gamma (z parametrami a i \(\lambda\)).

\subsection*{Rozwiązanie}
W tym zadaniu nie dzieje się absolutnie nic ciekawego, po prostu przepisujemy wzór na gęstość funkcji o danym rozkładzie, a potem wstawiamy \(e^{tx}\) do wzoru na \(E(x)\).\\
\textbf{Rozkład jednostajny}\\
Dla \(X \sim Uni(a,b)\) f. gęstości to \(f_X(z)=\frac{1}{b-a}\) dla \(z\in (a,b)\)
\[E(e^{tX})=\int_a^b e^{tx} \frac{1}{b-a} dx = \frac{1}{b-a}
	\int_a^b e^{tx} dx = \frac{1}{b-a} \left[\frac{e^{tx}}{t}\right]_a^b = \frac{e^{tb}-e^{ta}}{t(b-a)}\]
\textbf{Rozkład wykładniczy}\\
Dla \(X \sim Exp(\lambda)\) f. gęstości to \(f_X(z)=\lambda e^{-\lambda z}\) dla \(z>0\)
\[E(e^{tX})=\int_0^\infty e^{tx} \lambda e^{\lambda x} dx = \lambda\int_0^\infty e^{(t-\lambda)x} dx = \lambda \left[\frac{e^{(t-\lambda)x} }{t-\lambda }  \right]_0^\infty\]
Całka musi być skończona, ponieważ w p.p. wartość oczekiwana będzie nieokreślona. Musimy więc mieć \(t<\lambda\)\footnote{na wykładzie chyba nie było to wprost powiedziane, ale \(\lambda > 0\), gdyby tak nie było to gęstość - a co za tym idzie prawdopodobieństwo - mogłyby wyjść ujemne, a tego byśmy \textbf{bardzo} nie chcieli}
teraz widać, że z takim ograniczeniem \(e^{(t-\lambda)x}\) zbiega do 0.
\[ \lambda \left[\frac{e^{(t-\lambda)x} }{t-\lambda }  \right]_0^\infty= \lambda \left(0 - \frac{1}{t-\lambda}\right) = \frac{\lambda}{\lambda-t}\]
\textbf{Rozkład gamma}\\
Dla \(X \sim Gamma(a,\lambda)\) f. gęstości to \(f_X(z)=\frac{1}{\Gamma(a)} \lambda^a z^{a-1} e^{-\lambda z}\) dla \(z>0\) \\
\[E(e^{tX}) = \int_0^\infty e^{tx} \frac{\lambda^a}{\Gamma(a)} x^{a-1} e^{-\lambda x} dx = \frac{\lambda^a}{\Gamma(a)}
	\int_0^\infty e^{(t-\lambda)x} x^{a-1} dx\]
Funkcja jest trochę podobna do funkcji gamma. Żeby była jescze bardziej podobna wyłączamy \((\lambda-t)^a\) przed całkę.
\[\frac{\lambda^a}{\Gamma(a)}
	\int_0^\infty e^{(t-\lambda)x} x^{a-1} dx= \frac{\lambda^a}{\Gamma(a)(\lambda-t)^a}
	\int_0^\infty e^{-(\lambda-t)x} ((\lambda - t)x)^{a} \frac{dx}{x}=\frac{\lambda^a \Gamma(a)}{\Gamma(a)(\lambda-t)^a}=\left(\frac{\lambda}{\lambda-t}\right)^a\]

\subsection*{Odpowiedź}
Jednostajny: \(\frac{e^{tb}-e^{ta}}{t(b-a)}\)\\
Wykładniczy: \(\frac{\lambda}{\lambda-t}\)\\
Gamma: \(\left(\frac{\lambda}{\lambda-t}\right)^a\)
